\chapter{Methodology \& System Design}
\label{chapter:methodology}

This chapter presents the methodological framework and system design adopted for the development of a real-time AI-based simultaneous interpretation system. It formulates the problem from an engineering perspective as a low-latency, streaming speech processing task and describes the overall architecture used to transform continuous spoken input into translated speech output. The chapter details the design of the end-to-end interpretation pipeline, including audio acquisition, automatic speech recognition \cite{nayeem2025asr_survey}, stability-aware segmentation, neural machine translation \cite{stahlberg2020nmt_review}, and text-to-speech synthesis \cite{tan2021tts_survey}. Particular emphasis is placed on the modular and provider-agnostic design of the system, which enables flexible integration of multiple AI services while maintaining real-time performance constraints. Furthermore, the backend–client integration and communication mechanisms are described to illustrate how interpretation results are delivered to end users. Finally, the evaluation framework is introduced, outlining the metrics, experimental conditions, and ablation study design used to assess latency, stability, and robustness of the proposed system.

\section{Problem Definition and Design Objectives}
\label{sec:problem_definition}

The problem addressed is the design of a real-time AI-based simultaneous interpretation system that converts continuous spoken input in a source language into translated speech output in one or more target languages with minimal delay. From an engineering perspective, this task can be formulated as a streaming, low-latency signal processing and orchestration problem composed of multiple AI-based transformation stages operating on continuous input data.

Let $a(t)$ denote a continuous audio signal captured from a microphone, where $t$ represents time. The objective is to transform this signal into a translated audio output $\hat{a}_{L}(t)$ for a target language $L$, such that the perceived delay between the source speech and the translated output remains within acceptable bounds for live interpretation scenarios, typically on the order of one to three seconds, beyond which the interaction shifts toward consecutive interpretation rather than true simultaneity \cite{niehues2018simultaneousarxiv, ma2019staclarxiv}.


The system is structured as a sequential processing pipeline consisting of the following transformations:

\begin{enumerate}
    \item \textbf{Automatic Speech Recognition (ASR):}  
    The input audio stream $a(t)$ is converted into a stream of textual hypotheses:
    \[
    x(t) = \text{ASR}(a(t)),
    \]
    where $x(t)$ represents a sequence of partial and final text segments produced incrementally over time.

    \item \textbf{Segmentation:}  
    Due to the incremental nature of streaming ASR, the recognized text stream $x(t)$ contains unstable intermediate hypotheses. A segmentation function
    \[
    s_i = \text{SEG}(x(t))
    \]
    is therefore applied to extract stable and semantically coherent text segments $s_i$ suitable for translation.

    \item \textbf{Neural Machine Translation (NMT):}  
    Each extracted segment $s_i$ is translated into a target language $L$ using a neural translation model:
    \[
    y_i^{(L)} = \text{NMT}_L(s_i),
    \]
    where $y_i^{(L)}$ denotes the translated text segment.

    \item \textbf{Text-to-Speech (TTS):}  
    The translated text is synthesized into speech:
    \[
    \hat{a}_i^{(L)} = \text{TTS}_L(y_i^{(L)}),
    \]
    resulting in an audio signal that is delivered to the listener.
\end{enumerate}

The complete system can thus be viewed as a composite mapping:
\[
\hat{a}_L(t) = (\text{TTS}_L \circ \text{NMT}_L \circ \text{SEG} \circ \text{ASR})(a(t)).
\]

A key challenge in this formulation is that all transformations operate on \emph{streaming data} rather than fixed-length inputs. The ASR component produces partial hypotheses that may change over time, and each subsequent stage introduces additional latency. Consequently, the total end-to-end latency is the accumulated delay across all components, which must be minimized to ensure usability in live interpretation contexts \cite{ma2019staclarxiv}.

In addition to latency constraints, the system must support \emph{continuous operation} without relying on predefined sentence boundaries or offline processing. This requires careful coordination between asynchronous components and robust handling of partial and evolving intermediate results. Furthermore, the system should support \emph{multi-language output}, allowing the same recognized speech segment to be translated and synthesized into multiple target languages in parallel.

Based on this problem formulation and the main thesis objectives defined in Chapter~\ref{chapter:introduction}, the following \emph{system design objectives} were defined:

\begin{itemize}
    \item \textbf{Low end-to-end latency:} The system should minimize cumulative delays introduced by audio capture, recognition, segmentation, translation, and synthesis.
    \item \textbf{Streaming operation:} All components must operate incrementally on continuous input without requiring full utterances in advance.
    \item \textbf{Modular and extensible architecture:} Individual processing stages should be decoupled to allow independent replacement or extension.
    \item \textbf{Provider-agnostic AI components:} The system should support multiple ASR, NMT, and TTS providers through unified interfaces.
    \item \textbf{Multi-language output support:} Translation and synthesis should be performed concurrently for multiple target languages.
    \item \textbf{Hardware-compatible audio output:} Synthesized speech should be generated in standardized audio formats suitable for real-time playback and external routing.
\end{itemize}

To delimit the scope of this work, it is emphasized that this thesis does \emph{not} focus on training or optimizing machine learning models. Instead, it assumes the availability of pre-trained, cloud-based AI services and concentrates on the \emph{system-level design, orchestration, and real-time integration} of these components. The subsequent sections describe the architecture and implementation choices made to address the defined problem and achieve the stated design objectives.


\section{Overall System Architecture}
\label{sec:architecture}

This section presents the high-level architecture of the proposed real-time AI-based simultaneous interpretation system. It describes how the system is structured, how its main components interact, and how data flows through the pipeline under real-time constraints. The focus is on architectural organization, component separation, and concurrency rather than implementation details. This architectural perspective provides the foundation for understanding the design decisions that enable low-latency operation, modularity, and scalability in live interpretation scenarios.

\subsection{Architectural Overview and Component Structure}

The system is designed as a modular, real-time processing architecture that transforms continuous spoken input into translated speech output through a sequence of loosely cascaded components. Figure~\ref{fig:system_architecture} presents a high-level overview of the system architecture and illustrates the interaction between its main components.

\begin{figure}[!htbp]
  \centering
  \includesvg[width=\textwidth]{img/figures/system_architecture_diagram.drawio.svg}
  \caption[Real-time AI-based interpretation system architecture]{High-level architecture of the real-time AI-based simultaneous interpretation system.}
  \label{fig:system_architecture}
\end{figure}


At the core of the system is a streaming interpretation pipeline composed of five primary stages: audio acquisition, automatic speech recognition, segmentation, neural machine translation, and text-to-speech synthesis. Each stage is implemented as an independent subsystem with clearly defined responsibilities and interfaces, enabling separation of concerns and facilitating extensibility. This modular structure allows individual components to be replaced or extended without requiring changes to the overall system architecture \cite{Aminzadeh2009Advocate}.

The system operates on continuous audio input captured from a microphone, which is processed incrementally to satisfy real-time constraints. Rather than relying on offline or batch-based processing, all components are designed to operate on streaming data, enabling immediate propagation of information through the pipeline. This design choice is essential for minimizing end-to-end latency and ensuring responsiveness in live interpretation scenarios.

A backend orchestration layer coordinates the execution of the pipeline components and manages the distribution of intermediate and final results. This layer is responsible for maintaining the system lifecycle, routing data between components, and exposing the processed output to external clients \cite{he2019streamingarxiv}. Client applications, implemented as lightweight web-based interfaces, receive translated text and synthesized speech streams via network-based communication channels \cite{sathya2025realtime}.

This architecture emphasizes modularity, streaming operation, and real-time responsiveness. Detailed descriptions of the individual pipeline components are provided in Section~\ref{sec:pipeline}, while Figure~\ref{fig:system_architecture} serves as a structural reference for the system as a whole.


\subsection{Data Flow and Processing Stages}

The data flow within the system follows a sequential yet non-blocking processing model, as illustrated in Figure~\ref{fig:data_flow}. Live audio input enters the system as a continuous stream and propagates through a series of transformation stages that incrementally convert speech into translated audio output.

\begin{figure}[!htbp]
  \centering
  \includesvg[width=\textwidth]{img/figures/data_flow.drawio.svg}
  \caption[Interpretation pipeline Data flow ]{Data flow and concurrent execution model of the interpretation pipeline.}
  \label{fig:data_flow}
\end{figure}

Initially, the incoming audio stream is forwarded to the automatic speech recognition component, which continuously produces textual hypotheses. These hypotheses propagate to the segmentation stage, where stable text segments are extracted and forwarded downstream. The segmented text is then transmitted to the neural machine translation stage, which generates translated text for one or more target languages. Finally, the translated text is passed to the text-to-speech synthesis stage, producing audio output that is delivered to listeners.

The processing model combines streaming and event-driven behavior \cite{Sethiya2023EndtoEndSurvey}. While audio capture and speech recognition operate continuously on incoming data, downstream stages are activated by the availability of new text segments. This hybrid model enables timely propagation of information while avoiding unnecessary recomputation. Furthermore, translation and synthesis stages are executed independently for each target language, allowing multiple language outputs to be generated in parallel.

All processing stages communicate asynchronously, ensuring that no single component blocks the execution of others. This asynchronous data flow is critical for maintaining low latency and supporting concurrent processing across multiple languages. Figure~\ref{fig:data_flow} provides a conceptual view of this flow and highlights the propagation of data through the pipeline without detailing internal processing logic, which is addressed in subsequent sections.


\subsection{Architectural Design Rationale}
\label{sec:design_rationale}
The architectural decisions underlying the system design are guided by the requirements of real-time operation, extensibility, and deployment flexibility. A decoupled component structure was adopted to isolate responsibilities and reduce interdependencies between processing stages. Such modular architectures are widely used in real-time and distributed systems, as they simplify maintenance, enable independent evolution of components, and support the integration of alternative implementations without affecting the overall system behavior \cite{dragoni2017microservices, Aminzadeh2009Advocate, Sethiya2023EndtoEndSurvey}.

A streaming processing pipeline was selected to accommodate continuous speech input and minimize end-to-end latency. In contrast to batch-based approaches, streaming architectures allow incremental processing of incoming data and timely propagation of intermediate results. This design paradigm is commonly employed in real-time speech processing and simultaneous interpretation systems, where responsiveness is a critical usability factor \cite{fugen2007simultaneous, grissom2014realtime}.

To avoid reliance on a single technology provider, the system employs provider-agnostic interfaces for all AI-based components. Abstraction of external services is a common design practice in production AI systems, enabling flexibility, comparative evaluation, and resilience to changes in third-party services \cite{armbrust2020cloud}. This approach is particularly relevant in rapidly evolving AI ecosystems, where service capabilities and performance characteristics may vary.

Communication between the backend and client applications is implemented using WebSocket-based streaming. WebSockets\footnote{\url{https://datatracker.ietf.org/doc/html/rfc6455}} provide low-latency, bidirectional communication over persistent connections and are widely adopted in real-time web applications \cite{sathya2025realtime}. This communication model enables continuous delivery of both textual and audio data streams while maintaining a clear separation between processing logic and presentation.

Finally, the architecture is designed with hardware integration in mind. By standardizing audio formats and separating audio generation from playback and routing, the system aligns with common practices in professional audio and interpretation systems \cite{conference_interpreting_new_tech, wang2017realtimeaudio}. This design choice facilitates potential integration with external audio equipment such as mixers, interpretation consoles, or wireless transmission systems, which is further examined in Section~\ref{chapter:hardware}.


\section{Real-time Interpretation Pipeline}
\label{sec:pipeline}

This section presents the detailed design of the real-time interpretation pipeline proposed in this thesis. The pipeline follows a modular, streaming-oriented architecture in which audio input is incrementally processed through successive transformation stages, including speech recognition, segmentation, translation, and speech synthesis. Each stage is designed to operate asynchronously and under strict latency constraints, while maintaining robustness and compatibility with real-time deployment requirements. The following subsections describe each pipeline component in detail, highlighting its role, design rationale, and integration within the overall system.
 
\subsection{Audio Acquisition and Automatic Speech Recognition}

The first stage of the real-time interpretation pipeline is responsible for capturing live audio input and converting it into textual representations that can be processed by downstream components. In the proposed architecture, this stage combines audio acquisition and automatic speech recognition into a single coherent input module, reflecting their tight coupling in real-time speech processing systems.

%\paragraph{Audio Acquisition.}
Audio input is captured directly from a microphone device as a continuous stream. To ensure compatibility across the pipeline and with external speech recognition services, the audio signal is standardized to a single-channel (mono) pulse-code modulation (PCM) format with a sampling rate of 16~kHz maintaining a balanced compensation between audio quality and computational efficiency in the real-time speech system.

Rather than processing entire audio recordings, the input signal is segmented into small, fixed-duration 20ms chunks that are transmitted incrementally to the speech recognition component. This chunk-based streaming approach enables continuous processing of incoming speech and avoids the latency associated with batch-based or utterance-level audio handling. Audio acquisition is therefore designed to operate continuously and non-blockingly, ensuring that speech input is propagated through the system with minimal delay.

%\paragraph{Streaming Automatic Speech Recognition.}
The captured audio stream is forwarded to an automatic speech recognition (ASR) subsystem that operates in streaming mode. In contrast to offline recognition, streaming ASR produces textual output incrementally while the speaker is still talking \cite{weller2021streamingST}. This behavior is essential for simultaneous interpretation, as it allows downstream components to begin processing speech content before an utterance has fully completed. Figure~\ref{fig:asr_flow} illustrates the audio acquisition and streaming ASR stage, highlighting the incremental generation of recognition hypotheses and their propagation within the pipeline.

\begin{figure}[!htbp]
  \centering
  \includesvg[width=\textwidth]{img/figures/asr_flow.drawio.svg}
  \caption[Audio acquisition and streaming ASR pipeline]{Audio acquisition and streaming ASR pipeline}
  \label{fig:asr_flow}
\end{figure}

The ASR subsystem is integrated through a provider-agnostic abstraction layer, which exposes a unified interface for different speech recognition services. This abstraction decouples the system from any specific ASR provider and allows alternative implementations to be substituted without modifying the surrounding pipeline. Such flexibility is particularly important for evaluating robustness for multiple language translations under varying latency characteristics and for adapting to evolving service capabilities.

To evaluate robustness and flexibility under real-time conditions, the ASR subsystem integrates multiple cloud-based speech recognition services through a unified interface. The selected providers, summarized in Table~\ref{tab:asr_providers}, were chosen based on their support for streaming recognition, low-latency response, and availability of incremental transcription outputs. The use of multiple ASR providers allows the system to remain adaptable to different deployment constraints and service characteristics. This design choice further supports the provider-agnostic architecture described in Section~\ref{sec:design_rationale} and enables future substitution or extension without modifying downstream components.


\begin{table}[!htbp]
  \centering
  \caption[Integrated ASR providers and selection criteria]{Integrated automatic speech recognition providers and selection criteria}
  \label{tab:asr_providers}
  \begin{adjustbox}{max width=\textwidth}
    \begin{tabular}{p{4.5cm} p{11.8cm}}
      \toprule
      \textbf{ASR Provider} & \textbf{Selection Rationale} \\
      \midrule
      Google Speech-to-Text\footnotemark & Streaming recognition support and low-latency partial hypotheses\\
      \midrule
      Microsoft Azure Speech\footnotemark & Configurable streaming API and enterprise-grade deployment options\\
      \midrule
      AssemblyAI\footnotemark & Rich real-time transcription features and incremental output metadata \\
      \bottomrule
    \end{tabular}
  \end{adjustbox}
\end{table}

% Use the counter to decrement so the numbers match up correctly
\addtocounter{footnote}{-2}
\footnotetext{\url{https://cloud.google.com/speech-to-text/docs/streaming-recognize}}
\addtocounter{footnote}{1}
\footnotetext{\url{https://learn.microsoft.com/en-us/azure/ai-services/speech-service/}}
\addtocounter{footnote}{1}
\footnotetext{\url{https://www.assemblyai.com/docs/speech-to-text/streaming}}

Streaming ASR produces two types of outputs: \emph{partial hypotheses} and \emph{final hypotheses}. Partial hypotheses represent intermediate recognition results that may be revised as additional audio context becomes available, whereas final hypotheses correspond to stabilized recognition outputs that are no longer subject to change. As illustrated in Figure~\ref{fig:asr_flow}, partial hypotheses provide early access to speech content but are inherently unstable and unsuitable for direct translation.

%\paragraph{Role of Streaming in Real-Time Interpretation.}
Streaming ASR is a fundamental requirement for real-time interpretation systems due to strict latency constraints. Waiting for complete utterances before initiating recognition would introduce unacceptable delays between source speech and translated output. By processing audio incrementally and emitting partial hypotheses, the ASR subsystem enables early downstream processing while maintaining continuous system operation.

However, the use of streaming ASR introduces challenges related to output stability. Partial hypotheses may change significantly over time, leading to inconsistencies if translated prematurely \cite{grissom2014realtime}. As a result, the ASR output is not forwarded directly to the translation component. Instead, all recognition results are passed to a dedicated segmentation stage, which determines when textual segments are sufficiently stable and semantically coherent to trigger translation.

%\paragraph{Interface to Segmentation.}
The ASR subsystem therefore serves as the primary producer of textual input for the remainder of the interpretation pipeline. Both partial and final hypotheses are streamed to the segmentation component, along with metadata indicating hypothesis stability and timing. This separation of responsibilities allows the ASR stage to focus exclusively on speech-to-text conversion, while downstream components handle stability assessment and translation decisions. The segmentation strategy applied to these outputs is described in detail in Section~\ref{sec:segmentation}.



\subsection{Stability-Aware Segmentation Strategy}
\label{sec:segmentation}

In the proposed real-time interpretation pipeline, segmentation serves as the control mechanism that determines when recognized speech is forwarded for translation. Because the automatic speech recognition subsystem operates in streaming mode, it continuously emits partial hypotheses that may be revised as new audio arrives. Forwarding these hypotheses directly to the translation stage would result in unstable, duplicated, or semantically incomplete translations \cite{cohn2020perceptionTTS}. The segmentation strategy implemented in this system is therefore designed as an intermediate decision layer that explicitly controls translation triggering under real-time constraints.

%\paragraph{Segmentation as a Sequential Decision Process.}
Rather than applying independent segmentation rules in isolation, the system employs a sequential decision process in which multiple indicators are evaluated in a prioritized order. The primary objective is to identify speech segments that are sufficiently stable to preserve semantic meaning while remaining short enough to satisfy latency requirements. Figure~\ref{fig:segmentation_flow} provides an overview of this sequential segmentation process and the conditions evaluated at each decision stage.

\begin{figure}[!htbp]
  \centering
  \includesvg[width=\textwidth]{img/figures/segment_flow.drawio.svg}
  \caption[Segmentation sequential decision process]{Segmentation sequential decision process}
  \label{fig:segmentation_flow}
\end{figure}

At each update of the ASR output, newly recognized text is appended to a rolling buffer of unprocessed content. This buffer is continuously evaluated for potential segmentation points. The first level of segmentation candidates is derived from punctuation markers produced by the ASR system, which serve as initial indicators of clause boundaries. These punctuation cues provide a lightweight approximation of syntactic structure without relying on explicit linguistic parsing.

%\paragraph{Temporal and Clause-Level Constraints.}
Punctuation-based candidates are further refined using temporal information from the audio stream. Short pauses in speech are treated as reinforcement signals indicating the potential completion of a spoken unit. Segmentation is only considered when both textual cues (punctuation) and temporal cues (pauses) align, reducing the likelihood of premature segmentation at semantically incomplete positions, as illustrated in Figure~\ref{fig:segmentation_flow}.

To avoid excessive latency, the system enforces a maximum token length "K" that dynamically changes based on the speaker's speaking rate for any single segment. If the buffer exceeds this threshold without encountering a suitable boundary, a segmentation decision is forced at the nearest clause-level position. This mechanism ensures bounded delay while maintaining a preference for semantically coherent units over arbitrary time-based cuts.

In contrast, segments that fall below a minimum content threshold are discarded. Short interjections, fillers, and discourse markers that do not contribute meaningful semantic content are filtered prior to translation. This step reduces unnecessary translations and improves the fluency of the generated output.

%\paragraph{Stability Filtering and Translation Triggering.}
Once a candidate segment satisfies the structural and temporal constraints, a stability filtering stage is applied to determine whether the segment is sufficiently reliable for downstream processing. The system observes the candidate segment across successive updates of the streaming ASR output and verifies that its textual content remains unchanged over a predefined stability window. In the baseline configuration, this window corresponds to a minimum of four consecutive words for speakers with a normalized speaking rate in the range of $0.9\times$ to $1.1\times$ relative to nominal speech tempo. In Chapter~\ref{chapter:results}, an ablation study investigates the effect of changing the stability window length on the stability-latency trade-off. In the proposed system, To accommodate variations in speaking rate, the duration of the stability window is dynamically adjusted based on estimated speech tempo, thereby preventing unnecessary latency accumulation for slower speakers while maintaining robustness against premature finalization for faster speech. Only segments that remain stable throughout the applicable stability window are marked as finalized and forwarded to the translation stage. This final decision step is illustrated in Figure~\ref{fig:segmentation_flow} as the transition from candidate evaluation to segment finalization.


This stability filtering mechanism is central to preventing re-translation and incoherence. By delaying translation until the ASR output has converged, the system avoids forwarding text that is likely to be revised. At the same time, the stability window is kept intentionally short to avoid introducing unnecessary delay, thereby balancing responsiveness and reliability.

%\paragraph{Latency--Semantic Trade-Off Management.}
The segmentation strategy explicitly manages the trade-off between latency and semantic completeness. Translating complete sentences would maximize stability but introduce unacceptable delays in live interpretation scenarios. Translating individual words or unstable fragments would minimize latency but severely degrade coherence. The proposed approach resolves this trade-off by operating at the clause level, allowing translation to proceed incrementally while preserving meaning.

By prioritizing punctuation and pause-aligned boundaries, enforcing length constraints, and applying stability filtering, the system achieves early translation without sacrificing semantic integrity. Each segmentation decision reflects a compromise between immediacy and correctness, tailored to real-time interpretation requirements.

%\paragraph{Monotonic Progression and System Integration.}
Once a segment has been finalized and forwarded for translation, it is removed from the buffer and excluded from future segmentation decisions. This ensures a monotonic progression of processed content and prevents overlapping or duplicated translations. The segmentation component thus acts as a gatekeeper between recognition and translation, enforcing consistency across the pipeline.


\subsection{Neural Machine Translation}

The neural machine translation (NMT) stage is responsible for converting finalized speech segments produced by the segmentation component into translated text in one or more target languages \cite{stahlberg2020nmt_review}. In the proposed system, translation operates at the segment level rather than on complete sentences or continuous text streams. This design choice enables incremental translation while preserving semantic coherence and aligns with the real-time constraints of simultaneous interpretation.

%\paragraph{Segment-Level Translation.}
Each finalized segment emitted by the segmentation component is treated as an independent translation unit. Let $S_i$ denote a finalized speech segment in the source language. This segment is forwarded to the NMT subsystem, which produces one or more translated segments $\{T_i^{(L)}\}$, where $L$ denotes a target language. Segment-level translation allows downstream processing to begin immediately after stabilization, without waiting for complete utterances, thereby reducing end-to-end latency.
\[
\{T_i^{(L)}\} = \text{NMT}(S_i), \quad L \in \mathcal{L}
\]

\begin{figure}[!htbp]
  \centering
  \includesvg[width=\textwidth]{img/figures/nmt_flow.drawio.svg}
  \caption[Segment-level NMT fan-out with parallel target-language outputs]{Neural machine translation fan-out with parallel target-language outputs.}
  \label{fig:nmt_flow}
\end{figure}

Figure~\ref{fig:nmt_flow} illustrates the segment-level translation process, highlighting the fan-out of a single finalized segment into multiple target languages, the parallel execution of translation tasks, and the integration of glossary constraints within the translation flow.

%\paragraph{Asynchronous Translation Execution.}
Translation requests are executed asynchronously to avoid blocking the interpretation pipeline. Once a segment $S_i$ is emitted, translation tasks for all configured target languages are dispatched in parallel. This non-blocking execution model ensures that delays in translating into one language do not affect the availability of translations in other languages. As a result, the system maintains continuous throughput even under varying translation latencies, as reflected in the parallel execution paths shown in Figure~\ref{fig:nmt_flow}.

%\paragraph{Provider-Agnostic Translation Interface.}
The NMT subsystem is implemented using a provider-agnostic abstraction layer that exposes a unified translation interface independent of the underlying service implementation. This abstraction decouples the system from specific translation providers and allows alternative services to be integrated or substituted without modifications to the pipeline logic. The abstraction also enables comparative evaluation and flexible deployment across different environments.

%\paragraph{Selected NMT Providers and Rationale.}
Multiple cloud-based NMT services were integrated to ensure robustness, flexibility, and broad language coverage. The selected providers, summarized in Table~\ref{tab:nmt_providers}, were chosen based on their support for low-latency translation, multi-language coverage, and availability of programmatic interfaces suitable for real-time applications. Rather than optimizing for a single provider, the system design emphasizes adaptability and interoperability within a rapidly evolving AI service ecosystem.

\begin{table}[!htbp]
  \centering
  \caption[Integrated NMT providers and selection rationale]{Integrated neural machine translation providers and selection rationale}
  \label{tab:nmt_providers}
  \begin{adjustbox}{max width=\textwidth}
    \begin{tabular}{p{4.5cm} p{11.8cm}}
      \toprule
      \textbf{NMT Provider} & \textbf{Selection Rationale} \\
      \midrule
      Google Cloud Translation\footnotemark & Broad language coverage, low-latency API, glossary support \\
      \midrule
      Microsoft Azure Translator\footnotemark & Enterprise-grade deployment, configurable translation features \\
      \midrule
      DeepL API\footnotemark & High-quality translations for supported language pairs and terminology control\\
      \bottomrule
    \end{tabular}
  \end{adjustbox}
\end{table}

% --- Footnote Management ---
% Step 1: Move the counter back by (Number of Footnotes - 1)
\addtocounter{footnote}{-2} 
\footnotetext{\url{https://cloud.google.com/translate/docs}}

% Step 2: Increment and provide the next text
\stepcounter{footnote}
\footnotetext{\url{https://learn.microsoft.com/en-us/azure/ai-services/translator/}}

% Step 3: Increment and provide the final text
\stepcounter{footnote}
\footnotetext{\url{https://www.deepl.com/docs-api}}

%\paragraph{Multi-Language Fan-Out Strategy.}
To support simultaneous output in multiple target languages, the system employs a fan-out strategy at the translation stage. For each finalized segment $S_i$, translation requests are issued concurrently for all configured target languages. Each translated segment $T_i^{(L)}$ is processed independently and forwarded to the corresponding text-to-speech component. This design allows the system to scale naturally with the number of target languages and avoids serialization bottlenecks.

%\paragraph{Terminology and Glossary Integration.}
To improve translation consistency in domain-specific contexts, the proposed NMT subsystem integrates a zero-latency  glossary-based terminology mapping. Where predefined term source and target term mappings are applied pre-translation on the stable ASR output segments rendering technical terms, proper nouns, or organization-specific vocabulary. As illustrated in Figure~\ref{fig:nmt_flow}, glossary checks are applied within the translation flow without interfering with asynchronous execution, allowing terminology control without introducing additional latency.


\subsection{Text-to-Speech Synthesis and Audio Output}

The final stage of the real-time interpretation pipeline is responsible for converting translated text segments into audible speech and delivering them to the listener with minimal delay. This stage integrates neural text-to-speech (TTS) synthesis with real-time audio playback and routing, forming a unified output module that prioritizes low latency, and hardware compatibility \cite{tan2021tts_survey, conference_interpreting_new_tech}.

%\paragraph{Neural Text-to-Speech Abstraction.}
Translated text segments produced by the neural machine translation stage are forwarded to a neural text-to-speech subsystem. Similar to earlier pipeline components, TTS is accessed through a provider-agnostic abstraction layer that exposes a unified synthesis interface independent of the underlying service implementation. This design decouples the system from any specific TTS provider and allows alternative services to be integrated or replaced without modifying downstream audio handling logic.

%\paragraph{Selected TTS Providers and Rationale.}
Multiple cloud-based neural TTS services were integrated to ensure flexibility, voice diversity, and robust real-time performance. The selected providers, summarized in Table~\ref{tab:tts_providers}, were chosen based on their support for low-latency synthesis, multi-language voice availability, and programmatic access to raw audio outputs. Integrating multiple providers enables the system to adapt to different deployment requirements and service characteristics while maintaining a consistent synthesis interface.

\begin{table}[!htbp]
  \centering
  \caption[Integrated TTS providers and selection rationale]{Integrated text-to-speech providers and selection rationale}
  \label{tab:tts_providers}
  \begin{adjustbox}{max width=\textwidth}
    \begin{tabular}{p{4.5cm} p{11.8cm}}
      \toprule
      \textbf{TTS Provider} & \textbf{Selection Rationale} \\
      \midrule
      Microsoft Azure TTS\footnotemark & Low-latency neural voices and broad language support \\
      \midrule
      Google Cloud TTS\footnotemark & Stable real-time synthesis and multilingual voice availability \\
      \midrule
      ElevenLabs API\footnotemark & High-quality neural voices and expressive speech characteristics \\
      \bottomrule
    \end{tabular}
  \end{adjustbox}
\end{table}

% --- Manual Footnote Alignment ---
\addtocounter{footnote}{-2} % Adjust counter back to match the first mark
\footnotetext{\url{https://learn.microsoft.com/en-us/azure/ai-services/speech-service/text-to-speech}}

\stepcounter{footnote}
\footnotetext{\url{https://cloud.google.com/text-to-speech/docs}}

\stepcounter{footnote}
\footnotetext{\url{https://docs.elevenlabs.io}}

%\paragraph{Raw Audio Output and Voice Configuration.}
All TTS providers are configured to return raw pulse-code modulation (PCM) audio rather than compressed formats. This decision avoids decompression overhead and enables precise timing control during playback. As illustrated in Figure~\ref{fig:tts_flow}, audio is synthesized using a standardized sampling rate and channel configuration to ensure consistency across providers and compatibility with downstream routing components.

\begin{figure}[!htbp]
  \centering
  \includesvg[width=\textwidth]{img/figures/tts_flow.drawio.svg}
  \caption[TTS Synthesis and Audio Output Flow]{Text-to-Speech Synthesis and Audio Output Flow.}
  \label{fig:tts_flow}
\end{figure}

Voice selection is performed per target language, allowing different voices or speaking styles to be assigned depending on the output language. This mapping is resolved at runtime and remains independent of the synthesis provider, further reinforcing the modular design of the system.

%\paragraph{Latency Considerations in Speech Synthesis.}
Text-to-speech synthesis introduces additional latency due to network round-trip times, neural inference cost, and audio buffering. Several sources of latency were identified in the system, including acoustic look-ahead effects inherited from upstream components, sequential API calls between pipeline stages, neural speech generation time, and potential accumulation of audio packets during prolonged operation.

To mitigate these effects in the proposed system, synthesis requests are executed asynchronously and immediately forwarded to the audio playback subsystem upon completion. Persistent WebSocket connections are maintained to all cloud services to avoid repeated connection setup overhead. Additionally, synthesis parameters such as the speaking rate are dynamically adjusted with the speaking rate of the speakers (e.g., 1.1x–1.2x) to allow the playback pipeline to gradually reduce the accumulated buffer delay without compromising intelligibility.

%\paragraph{Real-Time Audio Playback and Buffer Management.}
Synthesized audio segments are scheduled for playback using a high-resolution audio clock to ensure precise timing. Rather than playing audio immediately upon arrival, each segment is queued and scheduled explicitly to avoid overlaps or gaps. Short overlaps between consecutive audio segments are applied to enable gapless stitching and prevent audible discontinuities and long-term buffer growth above the standard acceptable bounds for live interpretation scenarios stated in Section~\ref{sec:problem_definition}

%\paragraph{Separation of Synthesis and Playback.}
A strict separation is maintained between speech synthesis and audio playback. The TTS subsystem is responsible only for generating audio samples, while a dedicated audio output module handles buffering, scheduling, and device routing. This separation simplifies system reasoning, allows independent optimization of synthesis and playback, and enables flexible integration with different client platforms.

%\paragraph{Audio Routing and Hardware Compatibility.}
The final audio output stage supports configurable routing to different output devices, including speakers, virtual audio cables, or external hardware interfaces. By standardizing audio formats and isolating routing logic, the system can interface with professional interpretation equipment such as mixers, wireless transmitters, or interpretation consoles, which are discussed further in Chapter~\ref{chapter:hardware}. This design choice ensures that the system remains suitable for both software-only deployments and hybrid setups involving dedicated audio hardware.



\section{Backend and Client System Integration}
\label{sec:backend-client}

This section describes how the real-time interpretation pipeline is exposed to end users through a backend service and a browser-based client. All backend components discussed in Sections~\ref{sec:architecture} \& \ref{sec:pipeline} acts as the central coordination layer that manages pipeline execution, streaming communication, and language-specific routing, while the client provides a lightweight interface for interacting with the system and receiving interpretation results in real time.

%\paragraph{Backend Service Role.}
The backend service is responsible for orchestrating the real-time interpretation pipeline and managing communication between processing components and connected clients to validate the deployment of the system in realistic real-time simultaneous interpretation scenarios. It initializes and maintains the lifecycle of the interpretation pipeline, including audio ingestion, segmentation, translation, and synthesis. As illustrated in Figure~\ref{fig:backend-client}, the backend routes intermediate and final outputs to the appropriate clients based on language configuration.

\begin{figure}[!htbp]
  \centering
  \includesvg[width=\textwidth]{img/figures/client_flow.drawio.svg}
  \caption[Backend–client system integration]{Backend–client system integration, illustrating pipeline orchestration, language-based routing, and WebSocket-based streaming communication.}
  \label{fig:backend-client}
\end{figure}

By centralizing pipeline control, the backend ensures consistent state management across all processing stages and enables multiple clients to connect simultaneously without duplicating computational workloads.

%\paragraph{Pipeline Lifecycle Management.}
Pipeline execution is managed dynamically by the backend in direct response to client interactions. When a user initiates an interpretation session via the client interface, the backend establishes a persistent WebSocket connection and configures the interpretation pipeline according to the selected source and target languages. This connection serves as the primary communication channel through which audio input is streamed from the backend to the client and through which partial and finalized outputs are returned in real time on the client side. As the session progresses, the backend continuously streams recognition hypotheses, translated text segments, and synthesized audio to the client without requiring repeated connection setup. The pipeline remains active for the duration of the session and is gracefully terminated once the user ends it respectively, at which point associated resources are released. This WebSocket-based lifecycle management and persistent nature of WebSocket connections enables low-latency, bidirectional data exchange and supports scalable handling of multiple concurrent interpretation sessions \cite{sathya2025realtime}, in contrast to request–response protocols such as HTTP, which are ill-suited for continuous real-time streaming. A comparison between WebSockets and HTTP in the context of real-time interpretation systems is provided in Table~\ref{tab:websocket_vs_http}.

\begin{table}[!htbp]
  \centering
  \caption[Comparison of WebSockets \& HTTP for real-time communication]{Comparison of WebSockets and HTTP for Real-Time Interpretation Systems}
      \label{tab:websocket_vs_http}
  \begin{adjustbox}{max width=\textwidth}
    \begin{tabular}{p{6cm} p{4cm} p{6cm}}
      \toprule
      \textbf{Feature} & \textbf{WebSockets} & \textbf{HTTP} \\
      \midrule
      Latency & Low & Higher due to repeated requests \\
      \midrule
      Server push support & Native & Limited or emulated \\
      \midrule
      Streaming capability & Continuous & Fragmented \\
      \midrule
      Connection overhead & Single handshake & Repeated setup per request \\
      \midrule
      Suitability for real-time audio & High & Low \\
      \bottomrule
    \end{tabular}
  \end{adjustbox}
\end{table}


To support simultaneous interpretation into multiple target languages, the backend organizes communication using language-specific streaming channels. Each target language is associated with a dedicated WebSocket channel through which translated text and synthesized audio are transmitted. This channel-based organization simplifies client-side handling of multilingual output, enables parallel delivery of multiple languages, and allows clients to subscribe selectively to only the language streams relevant to a given session.

Client interaction with the system is provided through a browser-based interface that connects to the backend via persistent WebSocket streams. Through this interface, users can select the source and target languages, initiate and terminate interpretation sessions, and receive interpretation results in real time. Translated text segments are streamed incrementally as they become available, while synthesized speech is delivered continuously with minimal delay. This interaction model allows users to experience interpretation as a live, ongoing process rather than as a sequence of discrete requests.

A web-based client was chosen to maximize accessibility and deployment flexibility while maintaining low system complexity. Modern web browsers offer native support for real-time communication and audio playback, making it possible to deliver interpretation functionality without platform-specific installations. By confining all computationally intensive processing to the backend and limiting the client to interaction and presentation tasks, the system achieves a clear separation of concerns. This design facilitates maintainability, enables independent evolution of backend components, and supports consistent behavior across devices and operating systems.


\section{Evaluation Framework and Metrics Selection}
\label{sec:evaluation_framework}

This section defines the evaluation methodology and metrics used to assess the performance of the proposed real-time simultaneous interpretation system. The evaluation framework focuses on system-level behavior under real-time constraints rather than model-level accuracy. In particular, the evaluation emphasizes latency, stability, and robustness of the end-to-end pipeline and its individual components. The framework is designed to ensure that the reported results in Chapter~\ref{chapter:results} are reproducible, interpretable, and directly aligned with the system design objectives introduced in Chapter~\ref{chapter:methodology}.

\subsection{Evaluation Objectives}

The primary objective of the evaluation is to quantify how effectively the system performs real-time simultaneous interpretation under realistic operating conditions. Given the real-time nature of the task, the evaluation prioritizes latency-related metrics and pipeline stability over traditional model-centric accuracy measures.

Specifically, the evaluation aims to:
\begin{itemize}
  \item Measure the system's end-to-end latency from spoken input to translated audio output.
  \item Analyze latency contributions of individual pipeline stages, including speech recognition, segmentation, translation, and synthesis.
  \item Evaluate segmentation parameters impact on the trade-off between latency and stability.
\end{itemize}

Rather than evaluating linguistic accuracy on a fixed dataset, the evaluation focuses on system-level performance under standardized real-time conditions. All configurations were tested using identical audio input, system parameters, and measurement points, ensuring comparability across baseline and alternative components.

\subsection{Test Setup and Experimental Conditions}

\subsubsection{Baseline and Alternative Configurations.}
To ensure a clear and reproducible evaluation, a baseline configuration is defined to represent the default and recommended system setup. In addition, selected alternative configurations are evaluated by replacing individual pipeline components while keeping the remaining components unchanged. This approach isolates the impact of specific components without introducing combinatorial complexity. The evaluated configurations and their corresponding component selections are summarized in Table~\ref{tab:evaluation_configs}.

\begin{table}[!htbp]
  \centering
  \caption[Evaluation configurations and component selection]{Baseline and alternative evaluation configurations with selected components}
      \label{tab:evaluation_configs}
  \begin{adjustbox}{max width=\textwidth}
    \begin{tabular}{p{3cm} p{4cm} p{2cm} p{4cm}}
      \toprule
      \textbf{Configuration} & \textbf{ASR} & \textbf{NMT} & \textbf{TTS} \\
      \midrule
      Baseline & Azure Speech-to-Text & DeepL API & Azure Text-to-Speech \\
      \midrule
      Alternative ASR & Google Speech-to-Text & DeepL API & Azure Text-to-Speech \\
      \midrule
      Alternative TTS & Azure Speech-to-Text & DeepL API & Google Text-to-Speech \\
      \bottomrule
    \end{tabular}
  \end{adjustbox}
\end{table}

\subsubsection{Test Conditions.}
All experiments were conducted using the fully integrated real-time interpretation system described in Chapter~\ref{chapter:methodology}. Rather than evaluating linguistic accuracy on a fixed dataset, the evaluation focuses on system-level performance under standardized real-time operating conditions. To ensure consistency across test runs while preserving realistic streaming behavior, a scripted live speech protocol was adopted. A single speaker read the same predefined speech script, with a duration of approximately five minutes, at a natural speaking rate of roughly 130–160 words per minute (WPM). The full script is provided in Appendix~\ref{appendix:script}. Minor natural variations in speaking pace were allowed to reflect realistic interpretation scenarios without introducing uncontrolled variability.

Audio input was captured using a fixed recording setup in a controlled acoustic environment. All runs were performed inside a professional sound-isolated interpretation booth, minimizing background noise and external interference. This setup ensures that variations in system behavior are attributable to the processing pipeline rather than to environmental or acoustic factors. The same microphone, audio interface, operating system configuration, and audio parameters were used throughout all experiments to eliminate hardware-induced variability.

To further ensure fairness and reproducibility, all experiments were conducted using an identical hardware and network configuration. Network conditions corresponded to a stable wired broadband connection without artificial throttling or induced packet loss. Each experimental configuration, including the baseline and alternative component setups, was executed in three independent runs. Measured performance metrics were aggregated across repetitions to reduce the influence of transient fluctuations in network latency or cloud service response times. This controlled evaluation protocol ensures a fair and comparable assessment of system behavior across configurations.

A summary of the experimental setup and test conditions is provided in Table~\ref{tab:test_conditions}, which details the hardware environment, acoustic setting, audio parameters, network characteristics, language configuration, and segmentation parameters used throughout the evaluation.

\begin{table}[!htbp]
  \centering
  \caption[Test conditions and experimental setup]{Experimental setup and test conditions used for system evaluation}
    \label{tab:test_conditions}
  \begin{adjustbox}{max width=\textwidth}
    \begin{tabular}{p{5cm} p{7.3cm}}
      \toprule
      \textbf{Parameter} & \textbf{Value} \\
      \midrule
      CPU & Intel Core i7 (8 cores, 3.6 GHz) \\
      \midrule
      Microphone & Sennheiser Profile USB Microphone \\
      \midrule
      Testing Room & Audipack Silent 9700 Silent-Booth \\
      \midrule
      Operating System & Windows 11 (64-bit) \\
      \midrule
      Network Type & Wired broadband connection \\
      \midrule
      Average Network Latency & $\sim$20 ms (round-trip) \\
      \midrule
      Audio Sampling Rate & 16 kHz \\
      \midrule
      Audio Channels & Mono \\
      \midrule
      Audio Chunk Size & 20 ms \\
      \midrule
      Segmentation Stability Window & 4 words \\
      \midrule
      Source Language & English \\
      \midrule
      Target Languages & German, Arabic \\
      \midrule
      Session Duration (per run) & 5 minutes - with a 130-160 WPM speaking rate\\
      \midrule
      Runs Per Configuration & 3 Runs \\
      \bottomrule
    \end{tabular}
  \end{adjustbox}
\end{table}


\subsection{Evaluation Metrics}
\label{sec:evaluation_metrics}

The assessment of a cascaded speech-to-speech translation pipeline traditionally involves a modular approach, where the performance of each individual component is verified against static benchmarks. For the Automatic Speech Recognition (ASR) stage, accuracy is typically quantified using Word Error Rate (WER) \footnote{\url{https://www.speechmatics.com/company/articles-and-news/the-problem-with-word-error-rate-wer}} or Character Error Rate (CER) to measure the discrepancy between the hypothesis and a ground-truth transcript \cite{wer_definition_example}. Machine Translation (MT) quality is commonly evaluated using n-gram overlap metrics such as BLEU\cite{papineni2002bleu} and METEOR\cite{lavie2007meteor}, which assess linguistic fidelity and ordering. To ensure that the core intent remains intact across the pipeline, Semantic Fidelity is often measured through model-based approaches like BERTScore \cite{zhang2020bertscore}, which utilizes contextual embeddings to capture meaning beyond exact word matches. Finally, the Audio Quality of the synthesized output is evaluated using Mean Opinion Score (MOS) \cite{itutp8002} or PESQ \cite{rix2001perceptual} to determine the naturalness and intelligibility of the final speech signal.

However, in the context of simultaneous interpretation, these traditional model-centric metrics are insufficient as they assume an offline environment where the entire input is available before processing begins. Because our evaluation focuses on real-time operational performance and the critical latency-stability trade-off, we prioritize system-level metrics that quantify incremental behavior, pipeline fluidity, and end-to-end delay. Unlike batch metrics, system-level evaluation accounts for the temporal constraints where a correct translation delivered too late is considered a failure in a live scenario. The following system evaluation metrics, summarized in Table \ref{tab:system_metrics}, are drawn from the state-of-the-art literature in simultaneous speech translation to provide a holistic view of the pipeline's responsiveness.



\subsubsection{End-to-End System Latency}
\label{subsec:eval_e2e_latency}

The first evaluation objective focuses on quantifying the overall latency experienced by end users during real-time interpretation. End-to-end latency directly reflects how closely the translated output follows the source speech and is therefore a primary determinant of usability in simultaneous interpretation scenarios. To capture this behavior, latency metrics are selected that penalize excessive delay and reward sustained synchronization with the speaker.

A central metric used in this evaluation is \emph{Average Lagging (AL)}, which has become a standard measure in the simultaneous translation literature \cite{papi-etal-2022-laal}. AL quantifies the temporal delay of a system by comparing the emission time of target words against an idealized oracle translator operating at a constant rate. Formally, Average Lagging is defined as
\begin{equation}
\text{AL} = \frac{1}{\tau} \sum_{i=1}^{\tau} \left( t_i - \frac{i}{r} \right),
\end{equation}
where $t_i$ denotes the time at which the $i$-th target word is emitted, $r$ represents the average source-to-target token ratio, and $\tau$ is the cutoff index at which the system has consumed the entire source input. This metric is particularly relevant because it closely corresponds to the concept of \emph{Ear--Voice Span (EVS)} in human interpretation. In the proposed system, AL serves as an indicator of whether the generated interpretation maintains a sustainable delay that does not overload the listener’s working memory \cite{papi-etal-2022-laal}.

To address known limitations of Average Lagging, particularly its sensitivity to output length variation, \emph{Length-Adaptive Average Lagging (LAAL)} \cite{papi2022evaluation} is also employed. LAAL introduces a normalization factor based on the reference translation length, thereby correcting the tendency of AL to underestimate latency for overly verbose outputs. By accounting explicitly for translation length differences, LAAL enables fair comparison across configurations that may generate shorter or longer target sequences. This refinement is essential in real-time systems, where segment-level translation often leads to variable output lengths depending on segmentation and stabilization strategies \cite{ma2019staclarxiv, papi2022evaluation}.

In addition to lag-based measures, system efficiency is evaluated using the \emph{Real-Time Factor (RTF)} \cite{arriaga2025assessing}, which relates total processing time to the duration of the input audio. RTF is defined as
\begin{equation}
\text{RTF} = \frac{T_{\text{processing}}}{T_{\text{audio}}},
\end{equation}
where $T_{\text{processing}}$ denotes the cumulative processing time of the system and $T_{\text{audio}}$ represents the length of the input speech signal. An RTF value below unity indicates that the system can process audio faster than real time, whereas values greater than one imply that the system will progressively fall behind during extended sessions. In the context of live interpretation, maintaining $\text{RTF} < 1$ is a necessary condition for sustained usability.

\subsubsection{Pipeline Stage Latency Breakdown}
\label{sec:eval_stage_latency}

While end-to-end latency provides a holistic measure of system performance, it does not reveal how delays accumulate across individual pipeline stages. The second evaluation objective therefore decomposes latency into component-level contributions, allowing identification of bottlenecks within the automatic speech recognition, translation, and synthesis stages.

Responsiveness of upstream components is assessed using \emph{First-Token Latency (FTL)} \footnote{\url{https://docs.nvidia.com/nim/benchmarking/llm/latest/metrics.html}}, which measures the elapsed time between the onset of a source utterance and the availability of the first corresponding translated token \cite{chen2024core}. This metric captures how quickly the system reacts to incoming speech and is particularly important in simultaneous settings, where delayed initial output can significantly degrade user experience even if subsequent translation proceeds smoothly.

For the speech synthesis stage, latency is quantified using \emph{Time to First Audio (TTFA)}\footnote{\url{https://docs.nvidia.com/nim/benchmarking/llm/latest/metrics.html}}. TTFA measures the interval between the generation of a translated text segment by the machine translation component and the initiation of audible playback by the text-to-speech engine. Since the final system output is spoken audio, TTFA reflects the responsiveness of the output stage and directly influences the perceived continuity of interpretation \cite{chen2024core}. Measuring this delay enables targeted optimization of buffering strategies and audio scheduling mechanisms in the synthesis pipeline.

\subsection{Ablation Studies: Stability and Segmentation Behavior}
\label{sec:eval_stability}

In addition to baseline performance evaluation, ablation studies are conducted to analyze the sensitivity of the system to key segmentation parameters. These studies focus exclusively on parameters that influence the stability-latency trade-off and do not involve changes to model architectures or providers.

\subsubsection{Stability Window Duration Sensitivity}
As incremental ASR and translation often generate provisional outputs that may be revised as additional context becomes available, therefore, the stability window duration was introduced as in Section~\ref{sec:segmentation} to determine how long a candidate segment must remain unchanged before being forwarded for translation. In this ablation study the aim is to evaluate its impact, where multiple stability window durations are tested using the \emph{Flicker Rate}, a metric that captures how frequently provisional tokens are modified before being finalized \cite{rausei2023ralle}. High flicker rates correspond to frequent revisions in the output stream, which increase cognitive load and visual fatigue for users. In the proposed system, this metric is used to evaluate the effectiveness of segmentation and stabilization mechanisms and to identify operating points that balance responsiveness against excessive revisions.

Complementing flicker rate, \emph{Normalized Erasure (NE)} quantifies the amount of re-editing performed during incremental generation \cite{zhang2021beyond, arivazhagan2019retranslation}. NE measures the normalized edit distance between successive output states and is defined as
\begin{equation}
\text{NE} = \frac{\sum_{i} \text{edits}_i}{\sum_{i} \text{tokens}_i},
\end{equation}
where $\text{edits}_i$ represents the number of deletions or substitutions applied at update step $i$. NE serves as a counterbalance to latency-oriented metrics, as aggressive low-latency segmentation typically increases the amount of rework required downstream.


\subsubsection{Dynamic Token Threshold Sensitivity}
The dynamic token threshold limits the maximum length of a speech segment before forced segmentation occurs. To assess its influence, different maximum token thresholds are evaluated under identical test conditions using \emph{Consecutive Wait (CW)}, which measures how many source tokens the system consumes before emitting target output \cite{kano2023averaged, li2022system, simtrans_quality_latency_2022}. CW provides insight into whether the segmentation strategy is overly conservative, waiting for excessive context, or overly aggressive, emitting fragments prematurely \cite{simtrans_quality_latency_2022}. This ablation study assesses the effectiveness of token-based segmentation in bounding latency while preserving meaningful translation units.


\begin{table}[!htbp]
  \centering
  \caption[System performance metrics]{System performance metrics for real-time simultaneous interpretation}
  \label{tab:system_metrics}
  \small
  \begin{adjustbox}{max width=\textwidth}
    \begin{tabular}{p{4.5cm} p{4.2cm} p{7.6cm}}
      \toprule
      \textbf{Metric Category} & \textbf{Metric} & \textbf{Description} \\
      \midrule
      \multirow{4}{*}{\textbf{End-to-End Latency}}
        & Average Lagging (AL) 
        & Measures the average delay of the system relative to an ideal wait-0 interpreter. \\ \cline{2-3}
        & Length-Adaptive Average Lagging (LAAL) 
        & Extension of AL that normalizes latency with respect to reference output length. \\ \cline{2-3}
        & Average Token Delay (ATD) 
        & Average temporal delay of each emitted output token relative to its corresponding source token. \\ \cline{2-3}
        & Real-Time Factor (RTF) 
        & Ratio between total processing time and input audio duration, indicating real-time capability. \\
      \midrule
      \multirow{4}{*}{\textbf{Pipeline Stage Breakdown}}
        & First-Token Latency (FTL) 
        & Time elapsed from audio input to the emission of the first recognized or translated token. \\ \cline{2-3}
        & Decision Delay 
        & Duration for which the segmentation module waits before confirming a translation boundary. \\ \cline{2-3}
        & Inference Time 
        & Computation time required to process a single audio chunk or text segment. \\ \cline{2-3}
        & Time to First Audio (TTFA) 
        & Delay between translated text availability and the start of synthesized audio playback. \\
      \midrule
      \multirow{4}{*}{\textbf{Stability and Segmentation}}
        & Flicker (Instability) Rate 
        & Frequency with which provisional output tokens are revised before finalization. \\ \cline{2-3}
        & Normalized Erasure (NE) 
        & Proportion of deleted or rewritten tokens during incremental decoding. \\ \cline{2-3}
        & Consecutive Wait (CW) 
        & Maximum number of source tokens consumed before an output token is emitted. \\ \cline{2-3}
        & Meaning Unit (MU) Accuracy 
        & Degree to which emitted segments align with natural semantic or clause-level boundaries. \\
      \bottomrule
    \end{tabular}
  \end{adjustbox}
\end{table}


%%%%%%%%%%%%%%%%%%%%%%%%%%%%%%%%%%%%%%%%%%%%%%%%%%%%%
